\section{From six birds to agency}
\label{sec:dictionary}

\SBT\ \citep{six_birds_theory} explains macroscopic structure through six emergence mechanisms. Our approach to agency follows the same discipline: we assume only a finite stochastic substrate and ask which induced variables become stable, budgeted, and difference-making. This section states how each mechanism maps to a role in agency and how we probe that role in our experiments.

\paragraph{Theories and theory objects.}
In \SBT\ a \emph{theory} is a finite theory package: a carrier with a lens, a completion endomap, and an audit; its objects are the fixed points of the completion, the descriptions the theory recognises as complete at the chosen scale \citep{six_birds_theory}. In this paper a theory is the finite specialisation $T=(\Pi,L,\mathcal{F},B)$ of \cref{sec:engine}, and the packaging map of that section plays the role of the completion. Throughout the paper we reserve ``theory'' for the level $T$ and say \emph{theory object} for a stabilised package inside it. An \emph{agent} is a theory object with (i) an interface whose feasible settings are gated by a ledger and (ii) non-trivial difference-making within the level.

\paragraph{Three measurable aspects.}
We return repeatedly to three proxies.
\begin{itemize}[leftmargin=*, itemsep=0.25em]
  \item \textbf{Objecthood.} Does a lens that hides microscopic variables produce labels that behave like objects? We measure the idempotence defect of a packaging map $E$ (\cref{sec:engine,sec:ex_packaging}).
  \item \textbf{Maintained existence.} Is there a set of safe states that can be kept safe, with certainty under the model, using affordable actions? We measure the size of the viability kernel $|\K|$ (\cref{sec:engine,sec:ex_packaging,sec:ex_sweep}).
  \item \textbf{Difference-making.} Do affordable action sequences produce distinguishable outside futures? We measure feasible empowerment (\cref{sec:engine,sec:ex_holonomy,sec:ex_learning}).
\end{itemize}
The three are deliberately distinct. Labels can be object-like without any control, control without persistence is fragile, and, as \cref{sec:ex_packaging} shows, a zero packaging defect can occur without any maintenance at all.

\begin{table}[t]
\centering
\small
\begin{tabularx}{\linewidth}{@{}>{\raggedright\arraybackslash}p{0.17\linewidth} >{\raggedright\arraybackslash}X >{\raggedright\arraybackslash}X@{}}
\toprule
\textbf{Mechanism} & \textbf{Role in agency} & \textbf{How it appears in this paper} \\
\midrule
P$_1$ Operator rewriting &
Skill as a change of the effective law: the same interface moves become more reliable. &
A skill level $\theta$ lowers the slip probability. We compare median empowerment across $\theta$ (\cref{sec:ex_learning}). \\
\addlinespace
P$_2$ Feasibility constraints &
Constraints define what counts as an action: a command the ledger cannot pay for is not executed. &
An action is feasible only if its cost does not exceed the ledger. Viability uses stepwise feasibility; empowerment uses an initial-budget rule. The cost ablation sets movement costs to zero while keeping all actions (\cref{sec:ex_ablations}). \\
\addlinespace
P$_3$ Protocol cycles &
Order matters: moves that do not commute create outcomes visible only over several steps. &
With the protocol on, the step length depends on a staged phase $\phi$, so LEFT then RIGHT differs from RIGHT then LEFT. We compare empowerment across horizons (\cref{sec:ex_holonomy}). \\
\addlinespace
P$_4$ Identity and staging &
Discrete tokens and stages give stable reference frames and natural time scales. &
Only staging is exercised: the phase $\phi$ sets the natural packaging horizon $\tau=2$. The identity sector has a single value in every reported experiment, so identity is not tested here. \\
\addlinespace
P$_5$ Closure &
A closure turns a coarse description into stable, object-like labels and invariant sets. &
The viability kernel is a greatest fixed point (proved in Lean, \cref{app:lean_viability}); the packaging map tests whether coarse labels are fixed under composition. \\
\addlinespace
P$_6$ Resource accounting &
Persistence is paid for: resources are earned and spent to repair damage. &
A ledger $r$ with income and costs; a REPAIR action that resets damage. Funded repair makes a coherent viability kernel non-empty (\cref{sec:ex_packaging,sec:ex_sweep}). \\
\bottomrule
\end{tabularx}
\caption{The six mechanisms of \SBT, the role each plays in agency, and how this paper probes it. Coverage is uneven: P$_3$ and P$_6$ are tested by matched comparisons, P$_1$ in a stylised way, and P$_4$ only through phase staging.}
\label{tab:dictionary}
\end{table}

\paragraph{Guarding against false positives.}
Induced notions are only useful if they survive obvious baselines, a point already made in the \SBT\ programme \citep{six_birds_life}. We therefore include two null tests (\cref{sec:ex_nulls}): a system with a single action must have zero empowerment, and an external schedule must not be counted as a choice.
